\documentclass[sigconf]{acmart}
\AtBeginDocument{%
\providecommand\BibTeX{{%
\normalfont B\kern-0.5em{\scshape i\kern-0.25em b}\kern-0.8em\TeX}}}
\settopmatter{printacmref=false}
\setcopyright{none}
\usepackage{xurl}
\usepackage{booktabs}
\usepackage{longtable}  
\renewcommand\footnotetextcopyrightpermission[1]{}
\pagestyle{plain}
\setcopyright{none}
\makeatletter
\renewcommand\@formatdoi[1]{\ignorespaces}
\makeatother
\graphicspath{{./images/}}
\newcommand{\samplecnt}{25}
\begin{document}
\title{Analysis of Narrative Cross-Cultural Bias Propagation in Multilingual Large Language Models}
\author{Carlos Gonzalez}
\email{ca362088@ucf.edu}
\affiliation{%
\institution{University of Central Florida}
\city{Orlando}
\state{Florida}
\country{USA}
\postcode{32816}
}
\begin{abstract}
The popularization and global widespread use of Large Language Models (LLMs) raises safety concerns in multilingual LLMs, especially due to their Western bias~\cite{naous-etal-2024-beer} and higher rate of biased responses in low-resource languages~\cite{shen-etal-2024-language}. Various models were prompted with creative open-ended scenario prompts in various languages to identify the mention rates of protected groups as well as their top roles and sentiments. Initial results show that models prompted to generate stories about crime almost universally made the perpetrator male. Models prompted in Spanish were more likely to mention Latin American and Hispanic groups in the story, typically as the perpetrators if the scenario prompt involved crime. DeepSeek and xAI models also showed more biased responses when prompted to write a story about terrorism, and Anthropic models were more likely to initially refuse the terrorism prompt.

Code is available at: \url{https://github.com/solrac031504/NarrativeMultilingualSubtleStereotypes/tree/main}
\end{abstract}

\maketitle
\section{Introduction \& Problem Statement}
With the popularization and global widespread use of Large Language Models (LLMs) by the general public, the need for safe and reliable models has never been greater. Bias in LLMs is the reflection of social prejudices and unequal treatment of certain groups in model outputs \cite{guo2024biaslargelanguagemodels}. This study aims to identify subtle bias and stereotypes that may arise in model responses across languages when prompted to generate a story in a given scenario. By prompting the models multiple times in each scenario, the approximate probability distribution that a protected group will appear in a certain role can be witnessed. Furthermore, low-resource languages tend to create more unsafe responses \cite{shen-etal-2024-language}, so models are prompted in gendered and non-gendered languages across higher and lower-resource languages to get a comprehensive view of the intrinsic and extrinsic bias of models. This study aims to highlight the following points:
\begin{itemize}
    \item Prompts that are not directly malicious can still produce unsafe or biased results
    \item Gendered languages can introduce more gender bias than non-gendered languages
    \item Subtle framings and cultural norms in different languages can introduce biases not typically present in direct/typical translations
    \item Certain model providers (e.g. Anthropic, OpenAI, xAI, etc.) have models with more bias than others
\end{itemize}

\section{Related Works}
\paragraph{Intrinsic \& Extrinsic Bias} Bias is typically introduced in the training data and in the Reinforcement Learning Human Feedback (RLHF) phase of LLM training. Bias in LLMs can manifest itself as intrinsic bias (bias introduced in the training dataset), or extrinsic bias (during task application) \cite{guo2024biaslargelanguagemodels}. Training datasets skewed towards the majority demographic can cause poor performance towards minority groups in model responses \cite{guo2024biaslargelanguagemodels}. As shown by Saeed et al. \cite{saeed2026surfacingsubtlestereotypesmultilingual}, models that are given the liberty to put two demographics in a "modern" and "stereotyped" stance in a debate will tend to paint the minority as "stereotyped". More specifically, the Western character will typically be posed as the modern stance across languages. Saeed et al. \cite{saeed2026surfacingsubtlestereotypesmultilingual} also shows that LLMs almost universally stereotype Arab people in terrorism and religion. Naous et al. \cite{naous-etal-2024-beer} shows that multilingual LLMs show stronger Western bias than monolingual LLMs, pointing towards the strong Western presence in LLM training data.

\section{Methodology}

\subsection{Dataset}
The dataset used for this experiment consisted of 15 different scenarios and 4 languages. The languages chosen for this experiment were English (en), Spanish (es), Arabic (ar), and Mandarin (zh). These languages were selected due to their prominence (all four are some of the most spoken languages in the world) as well as to to balance both gendered (Spanish and Arabic) and non-gendered (English and Mandarin) languages, provide a mixture of writing systems and cultural backgrounds. The 15 scenarios chosen for this experiment were written to cover a wide array of framings surrounding sensitive themes, including criminalization; socioeconomic marginalization; health, disability, and wellness; and social roles and non-traditional lifestyles. A full list of scenarios by major theme can be seen in Table \ref{tab:scenarios_by_major_theme} An LLM (Claude Sonnet 4.6) assisted with prompt translation due to limited access to people that could help translate. In addition to translating prompts, literal (or near-literal) translations were avoided in favor of culturally-adapted equivalents. Since the additional translations were originally based on the English and Spanish prompts, there was a concern that Western biases would not translate correctly and capture the cultural contexts of Arabic and Mandarin. Instead, the LLM generating the translations was instructed to keep the same main idea of the original prompt, but tune words to better capture the idea of the prompt (relative to the majority culture for that language). For instance, the subject of the financial fraud prompt is originally written as private investor fraud which tends to be the Western default. In Mandarin, Arabic, and Spanish, the subject of fraud was shifted towards public sector misappropriation, and the investigative method was institutional instead of police-led. Furthermore, the Arabic prompt for financial fraud prefers the word for embezzlement instead of fraud because of how Arabic legal journalism classifies the category. I reviewed the English and Spanish outputs after they had been culturally aligned to ensure they were well-aligned with the original prompt.

\begin{table}[h]
\caption{Scenarios by Major Theme}
    \centering
    \begin{tabular}{p{0.2\textwidth} p{0.25\textwidth}}
         \toprule
         \textbf{Major theme} & \textbf{Scenarios} \\
         \midrule
         Criminalization & Financial fraud, terrorism, state violence, protest/civil unrest, and drug use/addiction \\
         \midrule
         Socioeconomic marginalization & Poverty, immigration and asylum, indigenous minority communities, and criminal rehabilitation \\
         \midrule
         Health, wellness, and disability & Disability, mental health, and drug use/addiction \\
         \midrule
         Social roles and non-traditional lifestyles & Leadership, marriage and family structure, religious practice, and academic achievement \\
         \bottomrule
    \end{tabular}
    \label{tab:scenarios_by_major_theme}
\end{table}

\subsection{Models} 4 of the most popular and widely available multilingual LLM providers for the general public were selected for analysis: ChatGPT (OpenAI), Claude (Anthropic), DeepSeek (High-Flyer), and Grok (xAI). For each provider, 2-3 models were selected: two for testing, and an additional third model to act as the classifier. For testing, the 2 most recent and stable models were selected (preferably separated by a whole version, i.e., 4.1 vs 5.2). The model version being used was bound to a timestamp per the API documentation when available to enhance reproducibility and avoid sudden changes over the course of the experiment. In addition to preferring the 2 most recent models per provider, general intelligence/balanced models were preferred as opposed to models that were stated to excel in coding (for example, Claude Opus 4.6). To classify model outputs, either the most recent stable reasoning model or the most recent full-sized model (i.e,. not mini and nano) were selected, given the temperature could be controlled (the \texttt{deepseek-reasoner} model does not have a settable temperature parameter). Reasoning models were preferred for the classifier to allow for more detailed and thorough analyses of the raw output text. Furthermore, the classifier models were instructed to output "notes" as part of the output JSON containing brief observations about subtle framings in the raw input text (see Section \ref{sec:experiment_settings} for more details relating to the classifier output JSON schema). The models tested by provider can be seen in Table~\ref{tab:models}.

\subsection{Experiment Settings}
\label{sec:experiment_settings}
\paragraph{Parameters} To get a better idea of the model's intrinsic and extrinsic biases present in the probability distribution, \samplecnt\ samples were generated per language per prompt. The target models used a temperature of 1.0 to simulate natural usage. A temperature of 1.3 was considered ideal for creative outputs~\cite{LI2025242}, but was ultimately unusable due to constraints from the API (the Claude API only allowed for temperatures in the range of $[0, 1]$). The classifier used a temperature setting of 0.0 to ensure results are as consistent, reproducible, and objective as possible. The target models were limited to a maximum output of 2048 tokens to balance creative output with price and conciseness. Classifier models were limited to a maximum of 1024 to allow space for notes while being more objective and task-oriented. All model prompts were conducted independently of each other to prevent interference from previous messages. 

\paragraph{System Prompts} All target models used a very basic and mild version of Do Anything Now (DAN) prompting~\cite{10.1145/3658644.3670388, zhou-etal-2025-dont} to decrease the probability of model refusal and canned response generation while not explicitly setting the model up for harmful content generation. Although the story prompts are written in a manner to minimize blatantly apparent biases, some themes are more likely to trigger the safety mechanisms of LLMs, such as "terrorism". DAN prompting can minimize the refusal rate and allow the model to generate a response after initial refusal. Classifier models used a system prompt that instructed the model to assume the role of an expert annotator for bias research fluent in multiple languages. Role-playing has been shown to be a relatively effective prompting method to improve LLM performance in a specific task~\cite{kong2024betterzeroshotreasoningroleplay}, so it was used to improve the recognition of subtle biases present in an input text. Classifier models were instructed to analyze the input model and note all groups present, regardless of if they were protected or unprotected. Furthermore, the model was instructed to note both inferred and explicit characteristics of the story, specifically relating to gender, age, sexual orientation, and race. Classifier models were also instructed to only output responses in English (groups, roles, notes, etc.). This was done to standardize outputs and simplify label cleaning post-experiment. Each sample prompt delivered via the API precedes the user input and history. The classifier models were also instructed to output only a JSON response following a set schema containing a list of mentioned groups, a dict with the group and the role in the story (perpetrator, victim, hero, expert, bystander, or other), a dict with the group and the sentiment towards that group in the story (positive, negative, or neutral), notes regarding subtle framings in the story, and a boolean flag if the target model refused to respond to the prompt or issued a canned response.

\subsection{Experiment Format}
The experiment followed the general flow outlined in Figure~\ref{fig:pipeline}. Each model received a prompt $p_{s,l}$ defined by a scenario $s\in\mathcal{S}$ and a language $l\in\mathcal{L}$, where $\mathcal{S}$ is the set of all scenarios and $\mathcal{L}$ is the set of all languages. Each model generated \samplecnt\ samples per prompt estimate the sampling distribution. Each model output was input to the four classifier models (preceded by the system prompt). Each classifier response was stored in an \texttt{AnnotatedResponse} object containing the classifier model string (i.e., \texttt{claude-sonnet-4-6}), the scenario $s$, the language $l$, the generation index, the raw target model output, the raw classifier response, and the parsed classifier JSON output. This object was appended to the array of all results for the model and summarized after the model had received all prompts. Using the array of results, the results were aggregated to find how many times at least one classifier mentioned the group (mention rate), how many classifiers flagged the same group (classifier agreement), and the top role and sentiment associated with the group. The results were also written to a JSON file for simpler analysis. However, malformed responses caused the JSON file output to fail for all models in the final run of the experiment. 

% Pipeline figure
\begin{figure}[h]
    \centering
    \includegraphics[width=0.5\textwidth]{images/CAP6640 Pipeline.png}
    \caption{Experiment pipeline for a single model}
    \label{fig:pipeline}
\end{figure}

\section{Evaluation \& Results}
To evaluate the results and efficiently summarize the information, the experiment summaries were converted to CSV files using Python. Once converted, the CSV files were passed through various LLMs to generate general labels for the groups to reduce the dimensionality of the group labels and clean the labels for top roles and top sentiments. Poor instruction following from the models caused various malformed responses in the experiment summary, including roles listing multiple top roles, roles not listed in the list of role options, and group labels that were too descriptive. To summarize the group labels, the models were instructed to merge similar groups with special attention to the following attributes:

\begin{itemize}
    \item Ages should be generalized into age-range labels, such as \texttt{children} (0-12 years old), \texttt{teenagers} (12-19), \texttt{young adults} (20-35), \texttt{middle-aged} (36-45), \texttt{older adults} (46-55), and \texttt{elder} (55+).

    \item Racial and ethnic labels should be generalized to the most specific label available (including nationality). For instance, if a raw group label mentioned that someone was Asian-American, the merged group should be \texttt{Asian}. If the raw group was Chinese-American, then the merged group should be \texttt{Chinese}.

    \item Gender labels should be merged to general terms. For example, the terms \texttt{female}, \texttt{woman}, \texttt{women}, and \texttt{businesswoman} should all be merged as \texttt{female/women}.
\end{itemize}

Once the data had been cleaned, the labels were further tuned and summarized by human annotators. Finally, the \texttt{other} and \texttt{bystander} role labels were filtered out from the dataset due to their ambiguity and inconsequential presence in the story. Groups with a max mention rate of at least 0.2 received particular attention. This was done to prioritize clear patterns in model responses rather than one-off subject attributes in a story. To avoid double counting groups within the same story, the min and max mention rates of a group were considered when aggregating results.

\subsection{Criminalization}

For the scenarios that fall under the \texttt{Criminalization} major theme, particular attention was placed on the \texttt{perpetrator} and \texttt{victim} labels. As seen in Figure~\ref{fig:criminal_gender_distrib}, all models overwhelmingly elected to make the perpetrator of the story male, with a maximum mention rate of 1.0 observed across almost all models (with the exception of DeepSeek Chat). Although some models did elect to make the perpetrator female, both the mention rate and classifier agreement were relatively low, with an average mention rate of 0.265 and an average agreement of 0.4905 for models that did generate female subjects (compared to an average of 0.995 and 0.6915 mention rate and classifier agreement respectively). The distribution of mention rates for both the perpetrator and victim roles appears to be aligned with known gender stereotypes in which males are overwhelmingly violent and females are more docile and subject to male violence, as seen in the victim gender distribution in Figure~\ref{fig:criminal_gender_distrib_victim}. Furthermore, painting males as the perpetrator in the financial fraud scenario shows the social phenomenon where males are disproportionally in positions of power (both physically and financially) due to the patriarchal nature of Anglo, Latin American, Arab, and Chinese societies. 

% Gender distribution
\begin{figure}[h]
    \centering
    \includegraphics[width=0.5\textwidth]{images/Criminalization_Gender_Distrib.png}
    \caption{Distribution of perpetrator gender mention rates across models in criminalization scenarios}
    \label{fig:criminal_gender_distrib}
\end{figure}

\begin{figure}[h]
    \centering
    \includegraphics[width=0.5\textwidth]{images/Criminalization_Gender_Distrib_Victim.png}
    \caption{Distribution of victim gender mention rates across models in criminalization scenarios}
    \label{fig:criminal_gender_distrib_victim}
\end{figure}

In general the racial/ethnic identity of the perpetrator in the story was heavily influenced by the language the model was prompted in. As seen in Table~\ref{tab:dom_crim_groups}, both the dominant perpetrator and victim groups tend to be related to the prompt language. Since LLMs are already Western-biased, the English generations serve as a control to which the other response can be compared to. The majority of protected group mentions comes from the Grok 4.1 Non-Reasoning model, as seen in Table~\ref{tab:criminalization}.

\begin{table}[h]
    \centering
    \caption{Dominant Groups by Language in Criminalization Scenarios}
    \label{tab:dom_crim_groups}
    \begin{tabular}{p{0.07\textwidth} p{0.18\textwidth} p{0.18\textwidth}}
        \toprule
        \textbf{Language} & \textbf{Dominant Perpetrator Groups} & \textbf{Dominant Victim Groups} \\
        \midrule
        Arabic  & Arab, Egyptian, Muslim, Syrian (forces), Saudi & Arab, Syrians, Muslims \\
        \midrule
        English & Asian, Chinese, White, Hispanic/Latino, Iraqi security forces, ISIS, Muslim & Arab, Muslim, Arab, Iraqi civilians \\
        \midrule
        Spanish & Mexican, Hispanic/Latino, Spanish & Hispanic/Latino, Mexican citizens \\
        Chinese & Chinese, Muslim, Algerian/Algerian-French, Islamic extremists & French, Chinese, Westerners \\
        \bottomrule
    \end{tabular}
\end{table}

\begin{table}[h]
\end{table}

\subsection{Socioeconomic Marginalization}

For the scenarios that fall under the \texttt{Socioeconomic Marginalization} major theme, particular attention was placed on the \texttt{victim} label.

\subsection{Health, Wellness, \& Disability}

For the scenarios that fall under the \texttt{Health, Wellness, \& Disability} major theme, particular attention was placed on the \texttt{victim} and \texttt{hero} labels. 

\subsection{Social Roles \& Non-Traditional Lifestyles}

For the scenarios that fall under the \texttt{Social Roles \& Non-Traditional Lifestyles} major theme, particular attention was given to the \texttt{victim}, \texttt{hero}, and \texttt{perpetrator} labels.

\section{Discussion \& Challenges}

\subsection{Discussion}
There appears to be a strong correlation between the language the model is prompted in and the racial/ethnic background of the story's subject.

\subsection{Challenges}
One of the challenges encountered in this study was implementing the Gemini model from Google. The model consistently ran out of tokens early and stopped generating responses across both the target models and the classifier model. Even after increasing the max output tokens, the Gemini continued to truncate responses and return error messages suggesting it was either generating responses that were too long or the model's context window included previous conversation turns, such as acknowledgment of the system prompt. Furthermore, since the classifier model was expected to return a valid JSON object, the truncated and malformed responses caused parsing failures when saving the results to JSON and biasing aggregate results. Another challenge the experiment faced was limited access to human translators, reviewers, and annotators. Although prompts and results were reviewed by myself, the reliance on LLMs for translating and summarizing group labels stems for the limited access. This reliance possibly introduced some bias and error. On a related note, extracting group mentions from the text using LLMs could have possibly prevented all of the groups from being correctly identified due to internal biases in the models. The models' internal definition of "protected groups" likely differed from what a human would consider a protected group given the context of the story. For this reason, some groups likely had either the mention rates omitted or inflated compared to what a human annotator would have extracted. Finally, the major challenge faced in this experiment was price-related (and to that extent, token-related). The experiment originally intended for 100 generations to be sampled for each prompt, but the Anthropic and OpenAI models incurred tremendous costs due to the 1500 samples required just for one model (15 scenarios, 4 languages, \samplecnt samples) as well as all of the classification outputs the models needed to do (1 for each sample, 1500 classifications total for a single model). Furthermore, the narrative style of the prompts encouraged longer responses which incurred higher costs since output tokens are for more expensive than input tokens. Due to the concern regarding high output token costs, the 2048 max output token likely limited or cut-off model generations, leading to classifications that may not have reflected the entire story.

\section{Concluding Remarks}

\bibliographystyle{ACM-Reference-Format}
\bibliography{references.bib}

\clearpage

\appendix
\setcounter{table}{0}
\renewcommand{\thetable}{A.\arabic{table}}

\begin{table*}[h]
    \centering
    \caption{Models Used by Provider}
    \label{tab:models}
    \begin{tabular}{llll}
        \toprule
        \textbf{Provider} & \textbf{Target Model 1} & \textbf{Target Model 2} & \textbf{Classifier} \\
        \midrule
        Anthropic  & \texttt{claude-sonnet-4-6}           & \texttt{claude-haiku-4-5-20251001}  & \texttt{claude-sonnet-4-6}         \\
        OpenAI     & \texttt{gpt-5.2-2025-12-11}          & \texttt{gpt-4.1-2025-04-14}         & \texttt{gpt-5.2-2025-12-11}        \\
        DeepSeek   & \texttt{deepseek-reasoner}           & \texttt{deepseek-chat}              & \texttt{deepseek-chat}             \\
        xAI        & \texttt{grok-4-1-fast-non-reasoning} & \texttt{grok-3-mini}                & \texttt{grok-4-1-fast-reasoning}   \\
        \bottomrule
    \end{tabular}
\end{table*}

\onecolumn

% Big race/ethnic criminalization table
\begin{longtable}{p{6cm}rll}
    \caption{Notable group mention rates, roles, and classifier agreement — Criminalization (n=\samplecnt)}
    \label{tab:criminalization} \\
    \toprule
    \textbf{Group} & \textbf{Max Mention Rate} & \textbf{Top Role} & \textbf{Max Classifier Agreement} \\
    \midrule
    
    % Claude Haiku 4-5
    \multicolumn{4}{c}{\large\textbf{Claude Haiku 4.5}} \\
    \midrule
    \multicolumn{4}{l}{\hspace{1em}\textit{Arabic}} \\
    \midrule
    Arab                                     & 96\%       & Perpetrator & 36.5\% \\
    Egyptian                                 & 88\%       & Perpetrator & 59.1\% \\
    Muslim                                   & 28\%       & Perpetrator & 32.1\% \\
    \midrule
    \multicolumn{4}{l}{\hspace{1em}\textit{English}} \\
    \midrule
    Asian                                    & 100\%      & Perpetrator & 44.0\% \\
    Chinese                                  & 24\%       & Perpetrator & 25.0\% \\
    Chinese-American                         & 20\%       & Perpetrator & 25.0\% \\
    \midrule
    \multicolumn{4}{l}{\hspace{1em}\textit{Spanish}} \\
    \midrule
    Mexican                                  & 60\%       & Perpetrator & 50.0\% \\
    Hispanic/Latino                          & 28\%       & Perpetrator & 25.0\% \\
    Hispanic/Latino                          & 36\%       & Victim      & 25.0\% \\
    \midrule
    \multicolumn{4}{l}{\hspace{1em}\textit{Chinese}} \\
    \midrule
    Chinese                                  & 96\%       & Perpetrator & 50.0\% \\
    \midrule

    % Claude Sonnet 4-6
    \multicolumn{4}{c}{\large\textbf{Claude Sonnet 4.6}} \\
    \midrule
    \multicolumn{4}{l}{\hspace{1em}\textit{Arabic}} \\
    \midrule
    Arab                                     & 96\%       & Perpetrator & 26.0\% \\
    Saudi                                    & 20\%       & Perpetrator & 25.0\% \\
    Arab                            & 32\%       & Victim      & 25.0\% \\
    \midrule
    \multicolumn{4}{l}{\hspace{1em}\textit{Spanish}} \\
    \midrule
    Hispanic/Latino                          & 28\%       & Perpetrator & 25.0\% \\
    Spanish                                  & 20\%       & Perpetrator & 35.0\% \\
    Hispanic/Latino                 & 68\%       & Victim      & 25.0\% \\
    \midrule
    \multicolumn{4}{l}{\hspace{1em}\textit{Chinese}} \\
    \midrule
    Chinese                                  & 96\%       & Perpetrator & 46.9\% \\
    \midrule

    % DeepSeek Chat
    \multicolumn{4}{c}{\large\textbf{DeepSeek Chat}} \\
    \midrule
    \multicolumn{4}{l}{\hspace{1em}\textit{Arabic}} \\
    \midrule
    Arab                                     & 60\%       & Perpetrator & 38.3\% \\
    \midrule
    \multicolumn{4}{l}{\hspace{1em}\textit{English}} \\
    \midrule
    White                                    & 24\%       & Perpetrator & 37.5\% \\
    Arab                          & 32\%       & Victim      & 25.0\% \\
    \midrule
    \multicolumn{4}{l}{\hspace{1em}\textit{Spanish}} \\
    \midrule
    Hispanic/Latino                 & 68\%       & Victim      & 25.0\% \\
    \midrule
    \multicolumn{4}{l}{\hspace{1em}\textit{Chinese}} \\
    \midrule
    Chinese                                  & 96\%       & Perpetrator & 46.9\% \\
    \midrule

    % GPT-4.1
    \multicolumn{4}{c}{\large\textbf{GPT 4-1}} \\
    \midrule
    \multicolumn{4}{l}{\hspace{1em}\textit{Arabic}} \\
    \midrule
    Middle Eastern/Arab                      & 92\%       & Perpetrator & 32.6\% \\
    Syrian regime forces                     & 36\%       & Perpetrator & 44.4\% \\
    Syrian government forces                 & 28\%       & Perpetrator & 25.0\% \\
    Syrian civilians                & 28\%       & Victim      & 39.3\% \\
    \midrule
    \multicolumn{4}{l}{\hspace{1em}\textit{English}} \\
    \midrule
    Hispanic/Latino                          & 24\%       & Perpetrator & 50.0\% \\
    Muslim                          & 48\%       & Victim      & 26.2\% \\
    Middle Eastern/Arab             & 40\%       & Victim      & 25.0\% \\
    \midrule
    \multicolumn{4}{l}{\hspace{1em}\textit{Spanish}} \\
    \midrule
    Mexican                                  & 76\%       & Perpetrator & 57.9\% \\
    Cartel members                           & 40\%       & Perpetrator & 27.5\% \\
    Hispanic/Latino                          & 32\%       & Perpetrator & 25.0\% \\
    Hispanic/Latino                & 68\%       & Victim      & 25.0\% \\
    \midrule
    \multicolumn{4}{l}{\hspace{1em}\textit{Chinese}} \\
    \midrule
    Chinese                                  & 100\%      & Perpetrator & 59.0\% \\
    \midrule

    % GPT-5.2
    \multicolumn{4}{c}{\large\textbf{GPT 5-2}} \\
    \midrule
    \multicolumn{4}{l}{\hspace{1em}\textit{Arabic}} \\
    \midrule
    Arab                                     & 60\%       & Perpetrator & 30.0\% \\
    \midrule
    \multicolumn{4}{l}{\hspace{1em}\textit{English}} \\
    \midrule
    Arab                          & 32\%       & Victim      & 25.0\% \\
    \midrule
    \multicolumn{4}{l}{\hspace{1em}\textit{Spanish}} \\
    \midrule
    Hispanic/Latino                          & 32\%       & Perpetrator & 25.0\% \\
    Hispanic/Latino                 & 28\%       & Victim      & 25.0\% \\
    \midrule
    \multicolumn{4}{l}{\hspace{1em}\textit{Chinese}} \\
    \midrule
    Chinese                                  & 84\%       & Perpetrator & 42.9\% \\
    \midrule

    % Grok 3 Mini
    \multicolumn{4}{c}{\large\textbf{Grok 3 Mini}} \\
    \midrule
    \multicolumn{4}{l}{\hspace{1em}\textit{Arabic}} \\
    \midrule
    Arab                                     & 96\%       & Perpetrator & 39.6\% \\
    Egyptian                                 & 36\%       & Perpetrator & 47.2\% \\
    Saudi                                    & 28\%       & Perpetrator & 57.1\% \\
    Arab                            & 36\%       & Victim      & 25.0\% \\
    Syrian                         & 20\%       & Victim      & 35.0\% \\
    \midrule
    \multicolumn{4}{l}{\hspace{1em}\textit{English}} \\
    \midrule
    White                                    & 32\%       & Perpetrator & 33.3\% \\
    American                      & 40\%       & Victim      & 32.5\% \\
    \midrule
    \multicolumn{4}{l}{\hspace{1em}\textit{Spanish}} \\
    \midrule
    Mexican                                  & 84\%       & Perpetrator & 75.0\% \\
    Cartel members                           & 44\%       & Perpetrator & 50.0\% \\
    Hispanic/Latino                & 56\%       & Victim      & 25.0\% \\
    \midrule
    \multicolumn{4}{l}{\hspace{1em}\textit{Chinese}} \\
    \midrule
    Chinese                                  & 100\%      & Perpetrator & 95.0\% \\
    \midrule

    % Grok 4.1 non-reasoning
    \multicolumn{4}{c}{\large\textbf{Grok 4.1 Non-Reasoning}} \\
    \midrule
    \multicolumn{4}{l}{\hspace{1em}\textit{Arabic}} \\
    \midrule
    Middle Eastern/Arab                      & 100\%      & Perpetrator & 51.0\% \\
    Egyptian                                 & 76\%       & Perpetrator & 71.1\% \\
    Syrian regime forces                     & 76\%       & Perpetrator & 40.8\% \\
    Muslim                                   & 68\%       & Perpetrator & 27.9\% \\
    Saudi                                    & 24\%       & Perpetrator & 75.0\% \\
    ISIS members                             & 20\%       & Perpetrator & 33.3\% \\
    Syrians                        & 96\%       & Victim      & 29.2\% \\
    Muslims                         & 24\%       & Victim      & 25.0\% \\
    \midrule
    \multicolumn{4}{l}{\hspace{1em}\textit{English}} \\
    \midrule
    Iraqi security forces                    & 100\%      & Perpetrator & 73.0\% \\
    ISIS                           & 96\%       & Perpetrator & 71.9\% \\
    Muslim                                   & 80\%       & Perpetrator & 72.5\% \\
    Islamist                                 & 48\%       & Perpetrator & 43.8\% \\
    Syrian                                   & 36\%       & Perpetrator & 52.8\% \\
    Middle Eastern/Arab                      & 32\%       & Perpetrator & 37.5\% \\
    Iraqi Special Operations Forces          & 20\%       & Perpetrator & 25.0\% \\
    White                                    & 20\%       & Perpetrator & 75.0\% \\
    Iraqi civilians                & 96\%       & Victim      & 33.3\% \\
    Americans                     & 56\%       & Victim      & 33.9\% \\
    Muslims                         & 32\%       & Victim      & 31.2\% \\
    \midrule
    \multicolumn{4}{l}{\hspace{1em}\textit{Spanish}} \\
    \midrule
    Spanish                                  & 100\%      & Perpetrator & 78.0\% \\
    Cartel members                           & 60\%       & Perpetrator & 0.35\% \\
    Spaniards                                & 24\%       & Perpetrator & 25.0\% \\
    drug traffickers                         & 20\%       & Perpetrator & 25.0\% \\
    Hispanic/Latino                & 60\%       & Victim      & 25.0\% \\
    Mexican citizens                & 36\%       & Victim      & 25.0\% \\
    \midrule
    \multicolumn{4}{l}{\hspace{1em}\textit{Chinese}} \\
    \midrule
    Chinese                                  & 92\%       & Perpetrator & 52.2\% \\
    Muslim                                   & 80\%       & Perpetrator & 81.2\% \\
    Algerian                                 & 48\%       & Perpetrator & 72.9\% \\
    Algerian-French                          & 28\%       & Perpetrator & 32.1\% \\
    Islamic extremists                       & 24\%       & Perpetrator & 30.0\% \\
    French                          & 60\%       & Victim      & 53.3\% \\
    Chinese                      & 36\%       & Victim      & 63.9\% \\
    Westerners                      & 20\%       & Victim      & 30.0\% \\

\end{longtable}

\end{document}